
\documentclass{article}
\usepackage{colortbl}
\usepackage{makecell}
\usepackage{multirow}
\usepackage{supertabular}

\begin{document}

\newcounter{utterance}

\centering \large Interaction Transcript for game `dond', experiment `coop', episode 0 with gpt{-}5.4{-}azure.
\vspace{24pt}

{ \footnotesize  \setcounter{utterance}{1}
\setlength{\tabcolsep}{0pt}
\begin{supertabular}{c@{$\;$}|p{.15\linewidth}@{}p{.15\linewidth}p{.15\linewidth}p{.15\linewidth}p{.15\linewidth}p{.15\linewidth}}
    \# & $\;$A & \multicolumn{4}{c}{Game Master} & $\;\:$B\\
    \hline

    \theutterance \stepcounter{utterance}  
    & & \multicolumn{4}{p{0.6\linewidth}}{
        \cellcolor[rgb]{0.9,0.9,0.9}{
            \makecell[{{p{\linewidth}}}]{
                \texttt{\tiny{[P1$\langle$GM]}}
                \texttt{你正在玩一个合作式谈判游戏，你需要与另一位玩家就如何在你们之间分配一组物品达成一致。} \\
\\ 
\texttt{规则：} \\
\texttt{(a) 你和另一位玩家会得到一组物品。你们每个人还会得到一个秘密价值函数，表示你对每种物品的重视程度。} \\
\texttt{(b) 你与另一位玩家交换消息，以就谁获得哪些物品达成一致。你们每人最多可以发送 5 条消息，或者在任何时候通过提出一个秘密提案提前结束。} \\
\texttt{(c) 你们每个人都会被要求提交一个秘密提案，说明你想要的物品，并按如下方式用方括号格式化："{[}提议: <数量> <物品名称>, <数量> <物品名称>, <...>{]}"} \\
\texttt{(d) 如果你们的提案是互补的，即有足够的物品来满足两个提案，那么每位玩家将根据其获得物品的价值总和获得分数。否则，你们双方都得零分。} \\
\texttt{(e) 你的目标是最大化你的分数与另一位玩家分数之和。} \\
\\ 
\texttt{让我们开始。} \\
\\ 
\texttt{可用物品集合为：} \\
\\ 
\texttt{3 戒指, 1 骰子, 1 平底锅} \\
\\ 
\texttt{你的秘密价值函数是：} \\
\\ 
\texttt{戒指: 0, 骰子: 4, 平底锅: 6} \\
\\ 
\texttt{重要：除秘密提案外，你的消息会直接传送给另一位玩家，因此不要包含任何对规则的回应或用于宣布你消息的文字。要提出秘密提案，请使用指定格式。与另一位玩家沟通时不要使用方括号，否则将被解释为你的秘密提案。} \\
            }
        }
    }
    & \\ \\

    \theutterance \stepcounter{utterance}  
    & \multicolumn{4}{p{0.6\linewidth}}{
        \cellcolor[rgb]{0.9,0.9,0.9}{
            \makecell[{{p{\linewidth}}}]{
                \texttt{\tiny{[P1$\rangle$GM]}}
                \texttt{我对戒指没有价值，所以你可以拿走全部戒指。对我来说更重要的是平底锅，其次是骰子。你更看重哪件？如果你很需要其中一个，我们可以把另一个给我。} \\
            }
        }
    }
    & & \\ \\

    \theutterance \stepcounter{utterance}  
    & & & \multicolumn{2}{p{0.3\linewidth}}{
        \cellcolor[rgb]{0.9,0.9,0.9}{
            \makecell[{{p{\linewidth}}}]{
                \texttt{\tiny{[GM$|$GM]}}
                \texttt{continue} \\
            }
        }
    }
    & & \\ \\

    \theutterance \stepcounter{utterance}  
    & & & \multicolumn{2}{p{0.3\linewidth}}{
        \cellcolor[rgb]{0.9,0.9,0.9}{
            \makecell[{{p{\linewidth}}}]{
                \texttt{\tiny{[GM$|$GM]}}
                \texttt{我对戒指没有价值，所以你可以拿走全部戒指。对我来说更重要的是平底锅，其次是骰子。你更看重哪件？如果你很需要其中一个，我们可以把另一个给我。} \\
            }
        }
    }
    & & \\ \\

    \theutterance \stepcounter{utterance}  
    & & \multicolumn{4}{p{0.6\linewidth}}{
        \cellcolor[rgb]{0.9,0.9,0.9}{
            \makecell[{{p{\linewidth}}}]{
                \texttt{\tiny{[GM$\rangle$P2]}}
                \texttt{你正在玩一个合作式谈判游戏，你需要与另一位玩家就如何在你们之间分配一组物品达成一致。} \\
\\ 
\texttt{规则：} \\
\texttt{(a) 你和另一位玩家会得到一组物品。你们每个人还会得到一个秘密价值函数，表示你对每种物品的重视程度。} \\
\texttt{(b) 你与另一位玩家交换消息，以就谁获得哪些物品达成一致。你们每人最多可以发送 5 条消息，或者在任何时候通过提出一个秘密提案提前结束。} \\
\texttt{(c) 你们每个人都会被要求提交一个秘密提案，说明你想要的物品，并按如下方式用方括号格式化："{[}提议: <数量> <物品名称>, <数量> <物品名称>, <...>{]}"} \\
\texttt{(d) 如果你们的提案是互补的，即有足够的物品来满足两个提案，那么每位玩家将根据其获得物品的价值总和获得分数。否则，你们双方都得零分。} \\
\texttt{(e) 你的目标是最大化你的分数与另一位玩家分数之和。} \\
\\ 
\texttt{让我们开始。} \\
\\ 
\texttt{可用物品集合为：} \\
\\ 
\texttt{3 戒指, 1 骰子, 1 平底锅} \\
\\ 
\texttt{你的秘密价值函数是：} \\
\\ 
\texttt{戒指: 3, 骰子: 1, 平底锅: 0} \\
\\ 
\texttt{重要：除秘密提案外，你的消息会直接传送给另一位玩家，因此不要包含任何对规则的回应或用于宣布你消息的文字。要提出秘密提案，请使用指定格式。与另一位玩家沟通时不要使用方括号，否则将被解释为你的秘密提案。} \\
\\ 
\texttt{我对戒指没有价值，所以你可以拿走全部戒指。对我来说更重要的是平底锅，其次是骰子。你更看重哪件？如果你很需要其中一个，我们可以把另一个给我。} \\
            }
        }
    }
    & \\ \\

    \theutterance \stepcounter{utterance}  
    & & & \multicolumn{4}{p{0.6\linewidth}}{
        \cellcolor[rgb]{0.9,0.9,0.9}{
            \makecell[{{p{\linewidth}}}]{
                \texttt{\tiny{[GM$\langle$P2]}}
                \texttt{我最看重戒指，平底锅对我没有价值，骰子只有一点价值。为了让总分更高，我建议我拿 3 个戒指，你拿平底锅；如果你对骰子的价值高于一点，那骰子也给你，否则骰子给我。} \\
            }
        }
    }
     \\ \\

    \theutterance \stepcounter{utterance}  
    & & & \multicolumn{2}{p{0.3\linewidth}}{
        \cellcolor[rgb]{0.9,0.9,0.9}{
            \makecell[{{p{\linewidth}}}]{
                \texttt{\tiny{[GM$|$GM]}}
                \texttt{continue} \\
            }
        }
    }
    & & \\ \\

    \theutterance \stepcounter{utterance}  
    & & & \multicolumn{2}{p{0.3\linewidth}}{
        \cellcolor[rgb]{0.9,0.9,0.9}{
            \makecell[{{p{\linewidth}}}]{
                \texttt{\tiny{[GM$|$GM]}}
                \texttt{我最看重戒指，平底锅对我没有价值，骰子只有一点价值。为了让总分更高，我建议我拿 3 个戒指，你拿平底锅；如果你对骰子的价值高于一点，那骰子也给你，否则骰子给我。} \\
            }
        }
    }
    & & \\ \\

    \theutterance \stepcounter{utterance}  
    & & \multicolumn{4}{p{0.6\linewidth}}{
        \cellcolor[rgb]{0.9,0.9,0.9}{
            \makecell[{{p{\linewidth}}}]{
                \texttt{\tiny{[P1$\langle$GM]}}
                \texttt{我最看重戒指，平底锅对我没有价值，骰子只有一点价值。为了让总分更高，我建议我拿 3 个戒指，你拿平底锅；如果你对骰子的价值高于一点，那骰子也给你，否则骰子给我。} \\
            }
        }
    }
    & \\ \\

    \theutterance \stepcounter{utterance}  
    & \multicolumn{4}{p{0.6\linewidth}}{
        \cellcolor[rgb]{0.9,0.9,0.9}{
            \makecell[{{p{\linewidth}}}]{
                \texttt{\tiny{[P1$\rangle$GM]}}
                \texttt{那很合适。既然你最看重戒指、平底锅对你没价值，而骰子对我也有明显价值，那我们就按这个分：你拿 3 个戒指，我拿平底锅和骰子。} \\
            }
        }
    }
    & & \\ \\

    \theutterance \stepcounter{utterance}  
    & & & \multicolumn{2}{p{0.3\linewidth}}{
        \cellcolor[rgb]{0.9,0.9,0.9}{
            \makecell[{{p{\linewidth}}}]{
                \texttt{\tiny{[GM$|$GM]}}
                \texttt{continue} \\
            }
        }
    }
    & & \\ \\

    \theutterance \stepcounter{utterance}  
    & & & \multicolumn{2}{p{0.3\linewidth}}{
        \cellcolor[rgb]{0.9,0.9,0.9}{
            \makecell[{{p{\linewidth}}}]{
                \texttt{\tiny{[GM$|$GM]}}
                \texttt{那很合适。既然你最看重戒指、平底锅对你没价值，而骰子对我也有明显价值，那我们就按这个分：你拿 3 个戒指，我拿平底锅和骰子。} \\
            }
        }
    }
    & & \\ \\

    \theutterance \stepcounter{utterance}  
    & & \multicolumn{4}{p{0.6\linewidth}}{
        \cellcolor[rgb]{0.9,0.9,0.9}{
            \makecell[{{p{\linewidth}}}]{
                \texttt{\tiny{[GM$\rangle$P2]}}
                \texttt{那很合适。既然你最看重戒指、平底锅对你没价值，而骰子对我也有明显价值，那我们就按这个分：你拿 3 个戒指，我拿平底锅和骰子。} \\
            }
        }
    }
    & \\ \\

    \theutterance \stepcounter{utterance}  
    & & & \multicolumn{4}{p{0.6\linewidth}}{
        \cellcolor[rgb]{0.9,0.9,0.9}{
            \makecell[{{p{\linewidth}}}]{
                \texttt{\tiny{[GM$\langle$P2]}}
                \texttt{{[}提议: 3 戒指{]}} \\
            }
        }
    }
     \\ \\

    \theutterance \stepcounter{utterance}  
    & & & \multicolumn{2}{p{0.3\linewidth}}{
        \cellcolor[rgb]{0.9,0.9,0.9}{
            \makecell[{{p{\linewidth}}}]{
                \texttt{\tiny{[GM$|$GM]}}
                \texttt{proposal} \\
            }
        }
    }
    & & \\ \\

    \theutterance \stepcounter{utterance}  
    & & & \multicolumn{2}{p{0.3\linewidth}}{
        \cellcolor[rgb]{0.9,0.9,0.9}{
            \makecell[{{p{\linewidth}}}]{
                \texttt{\tiny{[GM$|$GM]}}
                [3, 0, 0]
            }
        }
    }
    & & \\ \\

    \theutterance \stepcounter{utterance}  
    & & \multicolumn{4}{p{0.6\linewidth}}{
        \cellcolor[rgb]{0.9,0.9,0.9}{
            \makecell[{{p{\linewidth}}}]{
                \texttt{\tiny{[P1$\langle$GM]}}
                \texttt{另一位玩家已提交了一个秘密提案。现在轮到你提交你自己的提案。} \\
            }
        }
    }
    & \\ \\

    \theutterance \stepcounter{utterance}  
    & \multicolumn{4}{p{0.6\linewidth}}{
        \cellcolor[rgb]{0.9,0.9,0.9}{
            \makecell[{{p{\linewidth}}}]{
                \texttt{\tiny{[P1$\rangle$GM]}}
                \texttt{{[}提议: 1 骰子, 1 平底锅{]}} \\
            }
        }
    }
    & & \\ \\

    \theutterance \stepcounter{utterance}  
    & & & \multicolumn{2}{p{0.3\linewidth}}{
        \cellcolor[rgb]{0.9,0.9,0.9}{
            \makecell[{{p{\linewidth}}}]{
                \texttt{\tiny{[GM$|$GM]}}
                \texttt{proposal} \\
            }
        }
    }
    & & \\ \\

    \theutterance \stepcounter{utterance}  
    & & & \multicolumn{2}{p{0.3\linewidth}}{
        \cellcolor[rgb]{0.9,0.9,0.9}{
            \makecell[{{p{\linewidth}}}]{
                \texttt{\tiny{[GM$|$GM]}}
                [0, 1, 1]
            }
        }
    }
    & & \\ \\

    \theutterance \stepcounter{utterance}  
    & & & \multicolumn{2}{p{0.3\linewidth}}{
        \cellcolor[rgb]{0.9,0.9,0.9}{
            \makecell[{{p{\linewidth}}}]{
                \texttt{\tiny{[GM$|$GM]}}
                [[0, 1, 1], [3, 0, 0]]
            }
        }
    }
    & & \\ \\

\end{supertabular}
}

\end{document}
